\documentclass{article}
% Source: Topic_05_Teacher_Edition.pdf, PDF pages 64-75 (printed pages 261A-268B)
\usepackage{amsmath}
\usepackage{amssymb}
\usepackage{graphicx}
\usepackage{tikz}
\usepackage{pgfplots}
\pgfplotsset{compat=1.18}
\usepackage{booktabs}
\usepackage{xcolor}

\newenvironment{lesson-meta}{}{}        % lesson code, title, objective, EQ, vocab
\newenvironment{item-analysis}{}{}      % Savvas's Example × DOK table
\newenvironment{model-discuss}{}{}      % Launch opener
\newenvironment{concept-box}{}{}        % in-lesson concept reference
\newenvironment{concept-summary}{}{}    % end-of-lesson summary
\newenvironment{example}[1]{}{}         % #1 = example number
\newenvironment{tryit}[1]{}{}           % #1 = tryit number
\newenvironment{practice}[2]{}{}        % #1 = practice number, #2 = DOK
\newenvironment{te-addendum}[2]{}{}     % #1 = anchor example, #2 = type
\newcommand{\answer}[1]{}               % answer key, inside any env
\newcommand{\placeholder}[2]{}          % #1 = type, #2 = description
\newcommand{\uncertain}[2]{#1}          % #1 = best guess, #2 = alternatives
\newcommand{\te}[1]{}                   % teacher-only inline annotation

\begin{document}
% Source page 64: 261A Lesson Overview
\begin{lesson-meta}
lesson: 5-5
title: Function Operations
standards: none printed
practices: none printed
objective: I CAN perform operations on functions to answer real-world questions.

essential question: How do you combine, multiply, divide, and compose functions, and how do you find the domain of the resulting function?

\textbf{VOCABULARY}
\item composite function
\item composition of functions
\textbf{TOPIC VOCABULARY}
\te{No topic vocabulary list is printed on these lesson pages.}
\begin{tabular}{ll}
Lesson Code & 5-5 \\
Title & Function Operations \\
Standards & none printed \\
I CAN Objective & I CAN perform operations on functions to answer real-world questions. \\
Essential Question & How do you combine, multiply, divide, and compose functions, and how do you find the domain of the resulting function? \\
Lesson Vocabulary & composite function; composition of functions \\
Topic Vocabulary & none printed \\
\end{tabular}
\end{lesson-meta}
\begin{te-addendum}{0}{GeneralizeETP}
\textbf{Lesson Overview: FOCUS}
Objective. Students will be able to:
\item Combine functions by addition, subtraction, multiplication, or division and identify the domain of the result.
\item Compose functions, specifying the order in which the functions are applied and describing the domain of the composite function.
Essential Understanding. Functions can be combined by operations ((+,-,\times,\div)) and by composition. The result of the operation or composition can be described as a single function. The domain of the result may be different from the domains of the original functions.
\textbf{COHERENCE}
Previously in this course, students:
\item Evaluated individual functions for numerical inputs.
\item Described the domains (and ranges) of individual functions.
In this lesson, students:
\item Combine functions by operations ((+,-,\times,\div)) and compositions and describe the domain of the resulting function.
Increasing Coherence. Examples 4, 5, and 6 are not necessarily expected to be addressed on high-stakes assessments.
Later in this topic, students will:
\item Analyze composed functions, including identities and inverses of compositions of functions.
\textbf{RIGOR}
This lesson emphasizes a blend of conceptual understanding and application.
\item Students develop and use notation for applying operations to combine functions and to compose functions. They also explain when and why the domain of a combined or composite function is not the same as the domains of the original functions.
\item Students apply combinations and composites of functions to model situations such as discounts and coupons.
\end{te-addendum}
\begin{te-addendum}{0}{ELLAddendum}
\textbf{Vocabulary Builder}
\te{Glossary is available at SavvasRealize.com. Program Resources.}
\begin{tabular}{ll}
REVIEW VOCABULARY English & Spanish \\
domain & dominio \\
substitution & sustitución \\
NEW VOCABULARY & \\
composite function & funcion compuesta \\
composition of functions & composición de funciones \\
\end{tabular}
\textbf{VOCABULARY ACTIVITY}
Have students review aspects of substitution to prepare them for learning the composition of functions. To start, ask them to discuss how they can use substitution to check a solution to an equation.
\answer{Replace the variable and see if the result is a true statement.}
Then help students extend that idea of substitution by showing them the function (f(x)=(x+1)^2) and asking them to explain how to replace (x) with each of the following.
1. (3a) \answer{(f(3a)=(3a+1)^2=9a^2+6a+1)}
2. (r+s) \answer{(f(r+s)=(r+s+1)^2=r^2+s^2+2rs+2r+2s+1)}
\end{te-addendum}
\begin{te-addendum}{0}{LearnTogether}
\textbf{Student Companion}
Students can do their work for the lesson in their Student Companion or on Savvas Realize.
% FIG 64 1000 738 1152 917
\placeholder{illustration}{Student Companion cover showing purple and blue glowing sphere on black, enVision Algebra 2 and Savvas labels.}
\end{te-addendum}
\begin{te-addendum}{0}{GeneralizeETP}
\textbf{Mathematics Overview}
Students add, subtract, multiply, divide, and compose functions. With function composition, they also determine the correct order in which to perform the functions. Then students identify the domain of the resulting functions.
Students use verbal descriptions and algebraic representations to describe the addition, subtraction, multiplication, division, and composition of functions.
\textbf{Applying Math Practices: Attend to Precision}
Students identify and focus on the individual terms of pairs of functions. They do this when combining like terms and when multiplying each term of one function by each term of another function. Students apply this same focus when simplifying quotients and compositions of functions.
\textbf{Look For and Make Use of Structure}
Students use properties of operations on polynomials to combine and compose functions.
\end{te-addendum}
% Source page 65: 261B Model & Discuss
\begin{te-addendum}{0}{LearnTogether}
\textbf{STEP 1 Explore. MODEL \& DISCUSS. INSTRUCTIONAL FOCUS}
Students explore what it means to combine two functions. They see that the combination (by subtraction) has a practical meaning, and they express the combination of two functions as a single function.
\textbf{STUDENT COMPANION} Students can complete the Model \& Discuss activity in their Student Companion.
\end{te-addendum}
\begin{te-addendum}{0}{GeneralizeETP}
\textbf{Before: WHOLE CLASS. Implement Tasks that Promote Reasoning and Problem Solving}
Q: What do the terms revenue and cost mean?
\answer{Revenue is income; cost is an expense of the company.}
Q: How is profit related to revenue and cost?
\answer{Profit equals revenue minus cost.}
\textbf{During: SMALL GROUP. Support Productive Struggle in Learning Mathematics}
Q: What is the difference between a fixed cost and a variable cost?
\answer{A fixed cost does not change, while a variable cost changes based on the number of items involved.}
Q: Do both the revenue and the cost functions have both fixed and variable terms?
\answer{No, the cost function has a fixed term and a variable term, but the revenue function has only a variable term.}
\textbf{For Early Finishers}
Q: Last month, Grooming USA made a profit of \$300. What will their profit be if they double the number of pets they groom this month? How many pets do they have to groom this month to double last month's profit?
\answer{\$1,600; 160 pets}
Q: Does their profit double when the number of pets that are groomed doubles?
\answer{No; they only need to groom 30 more pets to double their profit. When they double the number of pets they groom, their profit is increased by more than 5 times.}
\textbf{After: WHOLE CLASS. Facilitate Meaningful Mathematical Discourse}
Q: How many functions were generated in this activity
\answer{two functions: one for revenue and one for cost}
Q: How can these two functions be combined to make a profit function?
\answer{by subtracting the function for cost from the function for revenue}
\end{te-addendum}
\begin{model-discuss}
\textbf{MODEL \& DISCUSS}
In business, the term profit is used to describe the difference between the money the business earns (revenue) and the money the business spends (cost).
% FIG 65 1000 641 1122 791
\placeholder{illustration}{Man brushes an orange bulldog on a grooming table; sign GROOMING, 25 dollars per pet.}
A. Let (x) represent the number of pets groomed in a month at Grooming USA. Define a revenue function for the business.
\answer{A. (r(x)=25x)}
B. Materials and labor for each pet groomed cost \$15. The business also has fixed costs of \$1,000 each month. Define a cost function for this business.
\answer{B. (c(x)=15x+1000)}
C. Last month, Grooming USA groomed 95 pets. Did they earn a profit? What would the profit be if the business groomed 110 pets in a month?
\answer{C. No; grooming only 95 pets results in a loss of \$50 for Grooming USA. Grooming 110 pets results in a profit of \$100.}
D. Generalize Explain your procedure for calculating the profit for Grooming USA. Suppose you wanted to calculate the profit for several different scenarios. How could you simplify your process?
\answer{D. To find the profit from 95 pets, I calculated (r(95)-c(95)); to find the profit from 110 pets, I calculated (r(110)-c(110)). If I needed to find the profit for several different values, it would make sense to find a function rule for profit and apply that simplified rule rather than applying two different functions and finding the difference of their values.}
\te{Digital variant of A: Grooming USA charges \$25 for every pet that is groomed. Let x represent the number of pets groomed in a month. Define a revenue function for the business. Banner: Critique \& Explain is available at SavvasRealize.com. Go back; Enter your answer; STUDENT EDITION; SAMPLE STUDENT WORK.}
\end{model-discuss}
\begin{te-addendum}{0}{HabitsOfMind}
\textbf{HABITS OF MIND. Use with MODEL \& DISCUSS}
Q: Make Sense and Persevere A business "breaks even" when its revenue equals its costs. How many pets would Grooming USA have to groom in order to break even?
\answer{100 pets}
\end{te-addendum}
% Source page 66: 261 Understand & Apply
\begin{te-addendum}{0}{LearnTogether}
\textbf{STEP 2 Understand \& Apply. INTRODUCE THE ESSENTIAL QUESTION. Establish Mathematics Goals to Focus Learning}
Use operations with integers to discuss how combining integers by addition, subtraction, and multiplication always results in an integer, while combining integers by division does not always result in an integer. Tell students to compare those ideas to the results they see in this lesson when they combine functions by addition, subtraction, multiplication, and division.
\end{te-addendum}
\begin{te-addendum}{1}{PurposefulQuestions}
\textbf{Add and Subtract Functions. Pose Purposeful Questions}
Q: How do you combine functions by addition or subtraction?
\answer{Combine the like terms of the two functions.}
Q: How do you know that the combined function has the same domain as the two original functions?
\answer{The possible elements for the combined function are the same as the possible elements for the original functions.}
Q: What new notation is used in the example and what does it mean?
\answer{((f+g)x) means (f(x)+g(x)), and ((f-g)x) means (f(x)-g(x)).}
\te{Printed notation omits parentheses around x.}
\end{te-addendum}
\begin{te-addendum}{4}{PurposefulQuestions}
\textbf{Additional Examples: ADDITIONAL EXAMPLE 4}
Let (f(x)=2x^2-3) and let (g(x)=4x-2). Find (f(g(x))) and find (g(f(x))).
(f(g(x))=2(4x-2)^2-3)
(=2(16x^2-16x-4)-3)
(=32x^2\ \ 32x\ \ 5)
(g(f(x))=4(2x^2-3)-2)
(=8x^2-12-2)
(=8x^2-14)
\answer{(f(g(x))=32x^2\ \ 32x\ \ 5); (g(f(x))=8x^2-14)}
\te{Printed first expansion has -4 instead of +4, and the final line has missing signs between terms; likely (32x^2-32x+5). Source PDF page 66. Digital labels: Go back; Essential Question; Solution.}
Example 4 Help students practice evaluating composite functions.
Q: What is the first step for each problem?
\answer{Replace the expression inside the outer parentheses.}
\end{te-addendum}
\begin{te-addendum}{6}{PurposefulQuestions}
\textbf{ADDITIONAL EXAMPLE 6}
A member currently has 75,000 awards points. In what order should the member take the bonuses?
Airline Awards Members! Get Both Bonuses! Add 10,000 award points! Add 5\% of your award points!
First take 10,000 points, then take 5\% bonus.
(10,000+75,000=85,000)
(1.05(85,000)=89,250)
Take a 5\% bonus, then add 10,000 points.
(1.05(75,000)=78,750)
(78,750+10,000=88,750)
The overall bonus is 500 points more if you add the 10,000 points first, and then take the 5\% bonus.
\answer{Add the 10,000 points first, and then take the 5\% bonus.}
Example 6 Help students see how the order of two functions can affect the final value.
Q: How do you decide which bonus to take first?
\answer{The greater total will occur when 5\% is taken on the additional 10,000 award points.}
Q: Why do the two orders have different outcomes?
\answer{Adding 5\% first yields less because you are adding 5\% of a smaller number.}
\te{Source PDF page 66. Additional Examples are available at SavvasRealize.com. Go back; Essential Question; Solution.}
\end{te-addendum}
\begin{example}{1}
\textbf{Add and Subtract Functions}
How do you define the sum, (f+g), and the difference, (f-g), of the functions (f(x)=3x+4) and (g(x)=x^2-5x+2)?
A. What is the sum of (f(x)=3x+4) and (g(x)=x^2-5x+2)?
To define the sum of two functions with known rules, add their rules.
((f+g)(x)=f(x)+g(x))
(=(3x+4)+(x^2-5x+2))
(=x^2+(3x-5x)+(4+2))
(=x^2-2x+6)
The domain of (f) is (\{x\mid x\text{ is a real number}\}).
The domain of (g) is (\{x\mid x\text{ is a real number}\}).
So the domain of (f+g) is (\{x\mid x\text{ is a real number}\}).
The sum of the two functions is (x^2-2x+6).
\te{COMMUNICATE PRECISELY Defining a function includes describing its domain. The domain of (f\pm g) is the intersection of the domains of f and g.}
\answer{A. (x^2-2x+6); domain: all real numbers.}
% Source page 67: 262 Example 1 continued
B. What is the difference of (f(x)=3x+4) and (g(x)=x^2-5x+2)?
To define the difference of two functions with known rules, subtract their rules.
((f-g)(x)=f(x)-g(x))
(=(3x+4)-(x^2-5x+2))
(=3x+4-x^2+5x-2)
(=-x^2+8x+2)
The domain of (f) is (\{x\mid x\text{ is a real number}\}). The domain of (g) is (\{x\mid x\text{ is a real number}\}). So the domain of (f-g) is (\{x\mid x\text{ is a real number}\}).
The difference of the two functions is (-x^2+8x+2).
\answer{B. (-x^2+8x+2); domain: all real numbers.}
\te{Examples are available at SavvasRealize.com. Activity.}
\end{example}
\begin{tryit}{1}
Let (f(x)=2x^2+7x-1) and (g(x)=3-2x). Identify rules for the following functions.
a. (f+g)
b. (f-g)
\answer{a. (f+g=2x^2+5x+2). b. (f-g=2x^2+9x-4).}
\end{tryit}
\begin{te-addendum}{1}{ElicitEvidence}
\textbf{Elicit and Use Evidence of Student Thinking}
Q: A student said that (f+g) and (f-g) are shorthand for longer expressions. Explain what the student means.
\answer{(f+g) is a shorter way to express ((2x^2+7x-1)+(3-2x)), and (f-g) is a shorter way to express ((2x^2+7x-1)-(3-2x)).}
\end{te-addendum}
\begin{te-addendum}{2}{PurposefulQuestions}
\textbf{Multiply Functions. Pose Purposeful Questions}
Q: How does the example involve multiplying functions?
\answer{The problem describes three functions---demand, price, and revenue---and revenue can be calculated by multiplying demand times price.}
Q: How do you multiply two functions?
\answer{Multiply each term of one function by each term of the other function and combine like terms}
Q: For the demand function (d(x)) and the price function (p(x)), is (d(x)\cdot p(x)) equal to (p(x)\cdot d(x))?
\answer{Yes; the products are equal because of the Commutative Property of Multiplication.}
\end{te-addendum}
\begin{te-addendum}{3}{RtIExtend}
\textbf{Extend Student Thinking. USE WITH EXAMPLE 3}
Have students define and state the domains of the following quotients when (f(x)=x^2+2x-3), (g(x)=x^2-1), and (h(x)=x^2+4x+3).
a. (\frac f{g+h})
\answer{(\frac{x^2+2x-3}{2x^2+4x+2}); The domain is the set of all real numbers x such that (x\neq-1).}
b. (\frac{f-h}g)
\answer{(\frac{-2x-6}{x^2-1}); The domain is the set of all real numbers x such that (x\neq-1) and (x\neq1).}
c. (\frac{h-g}{g-f})
\answer{(\frac{2x+2}{-x+1}); The domain is the set of all real numbers x such that (x\neq1).}
Q: How do you find the zeros of each denominator?
\answer{Set the denominator equal to zero, and solve for the variable.}
Q: How do the zeros of the denominator affect the domain of the function?
\answer{The zeros of the denominator are excluded from the domain of the function.}
\te{Source PDF page 67.}
\end{te-addendum}
\begin{example}{2}
\textbf{Multiply Functions. APPLICATION}
The demand d, in units sold, for a company's new brand of cell phone at price x, in dollars, is (d(x)=5,000-10x). The cost to manufacture the cell phone is (c(x)=2,200+80x). What is the company's expected profit from cell phone sales in terms of the price, x?
% FIG 67 801 532 1074 615
\placeholder{diagram}{Profit equals Price times Demand minus Cost. Four colored icons show a dollar sign and bar chart (Profit), a phone with price tag (Price), a circle of hands holding phones (Demand), and phone circuit board (Cost). Equals, multiplication and minus signs appear between icons.}
The company's profit will equal the price of its cell phones multiplied by the demand for its phones less the cost to manufacture.
Profit = price (\times) demand (-) cost
The demand is the function (d(x)). The price is the function (p(x)).
\begin{tabular}{lll}
 & & Domains \\
Price: & (p(x)=x) & (p(x):0\leq x) \\
Demand: & (d(x)=5,000-10x) & (d(x):0\leq x\leq500) \\
Cost: & (c(x)=2,200+80x) & (c(x):0\leq x\leq500) \\
Profit: & (P(x)=p(x)\cdot d(x)-c(x)) & (P(x):0\leq x\leq500) \\
 & (=x(5,000-10x)-(2,200+80x)) & \\
 & (=5,000x-10x^2-2,200-80x) & \\
 & (=-10x^2+4,920x-2,200) & \\
\end{tabular}
\te{USE STRUCTURE You can use the Associative and Commutative Properties to add and multiply functions, since these operations are based on addition and multiplication of real numbers. Callouts: Price, demand, and cost cannot be negative. Domain is the intersection of the domains of p, d, and c.}
% Source page 68: 263 Example 2 continued
The profit the company will earn in terms of the cell phone price x is represented by (P(x)=-10x^2+4,920x-2,200).
\answer{(P(x)=-10x^2+4,920x-2,200)}
\te{Examples are available at SavvasRealize.com. Activity.}
\end{example}
\begin{tryit}{2}
Suppose demand, d, for a company's product at cost, x, is predicted by the function (d(x)=-0.25x^2+1,000), and price, p, that the company can charge for the product is given by (p(x)=x+16). Find the company's revenue function. Revenue is demand times price.
\answer{(d\cdot p=-0.25x^3-4x^2+1,000x+16,000)}
\end{tryit}
\begin{te-addendum}{2}{ElicitEvidence}
\textbf{Elicit and Use Evidence of Student Thinking}
Q: Is it possible for x to have the value 64?
\answer{No; if (x=64), then (d(64)=-0.25(64)^2+1,000=-24), and (d(x)) cannot be less than 0.}
\end{te-addendum}
\begin{te-addendum}{3}{PurposefulQuestions}
\textbf{Divide Functions. Pose Purposeful Questions}
Q: How are the domains of (g) and (\frac fg) different?
\answer{g can have domain values that make the function equal zero, but those values cannot be in the domain of (\frac fg).}
\end{te-addendum}
\begin{example}{3}
\textbf{Divide Functions}
How do you define the quotient (\frac fg) of the functions (f(x)=x-7) and (g(x)=2x^2-13x-7)?
To define the quotient of two functions, take the quotient of their rules: ((\frac fg)(x)=\frac{f(x)}{g(x)})
((\frac fg)(x)=\frac{f(x)}{g(x)}=\frac{x-7}{2x^2-13x-7})
(=\frac{x-7}{(2x+1)(x-7)})
(=\frac1{2x+1})
The quotient of (\frac fg) is (\frac1{2x+1}). The domain of (\frac fg) is the set of all values for which f, g, and (\frac fg) are defined, so g cannot be 0. This is the set of all real numbers x such that (x\neq7) and (x\neq-\frac12).
\te{COMMON ERROR You may think that the domain of (\frac fg) is the set of real numbers since f and g are defined for all real numbers. However, (x\neq-\frac12) or (x\neq7) because g is now in the denominator. Printed or in callout; likely and.}
\answer{(\frac1{2x+1}); all real numbers x such that (x\neq7) and (x\neq-\frac12).}
\end{example}
\begin{tryit}{3}
Identify the rule and domain for (\frac fg) for each pair of functions.
a. (f(x)=x^2-3x-18), (g(x)=x+3)
b. (f(x)=x-3), (g(x)=x^2-x-6)
\answer{a. (\frac fg=x-6); The domain is all real numbers x such that (x\neq-3). b. (\frac fg=\frac1{x+2}); The domain is all real numbers x such that (x\neq3) and (x\neq-2).}
\end{tryit}
\begin{te-addendum}{3}{ElicitEvidence}
\textbf{Elicit and Use Evidence of Student Thinking}
Q: Explain how you would find the domain of (\frac gf).
\answer{First, identify the domains of g and f. Then find the values that make (f(x)) equal zero and eliminate them from the domain.}
\end{te-addendum}
\begin{te-addendum}{3}{HabitsOfMind}
\textbf{HABITS OF MIND. Use with EXAMPLES 1--3}
Q: Communicate Precisely How are the domains of (f+g), (f-g), (f\cdot g), and (\frac fg) related to the domains of f and g?
\answer{If f and g have the same domain, the domains of (f+g), (f-g), and (f\cdot g) are that same domain. The domain of (\frac fg) is that same domain, but without any values for which g is zero. If f and g have different domains (as in Example 2), the domain for (f+g), (f-g), and (f\cdot g) is the intersection of the two domains. The domain of (\frac fg) is further restricted by removing any values for which g is zero.}
\end{te-addendum}
\begin{te-addendum}{4}{GeneralizeETP}
\textbf{Compose Functions. Build Procedural Fluency From Conceptual Understanding}
Q: What is the difference between the two notations (f(x)\cdot g(x)) and (f(g(x)))?
\answer{(f(x)\cdot g(x)) is the product of two expressions, and (f(g(x))) uses the function (g(x)) as the variable for the function (f(x)).}
\end{te-addendum}
\begin{example}{4}
\textbf{Compose Functions. CONCEPTUAL UNDERSTANDING}
Let (f(x)=x^2) and let (g(x)=x+1). Investigate what happens when you apply the rule for g and then the rule for f to a number or variable.
A. Find the values of (f(g(0))), (f(g(3))), (f(g(-2))).
For each value of x, first apply the rule for g. Then apply the rule for f to the result.
\begin{tabular}{lll}
x & (g(x)) & (f(g(x))) \\
0 & (g(0)=0+1=1) & (f(g(0))=1^2=1) \\
3 & (g(3)=3+1=4) & (f(g(3))=4^2=16) \\
-2 & (g(-2)=-2+1=-1) & (f(g(-2))=(-1)^2=1) \\
\end{tabular}
\te{LOOK FOR RELATIONSHIPS When finding the rule for (f(g(x))), work from the inside parentheses to the outside ones. Notice that (g(x)) takes the place of the variable x in the function (f(x)).}
\answer{A. 1; 16; 1.}
% Source page 69: 264 Example 4 continued
B. Find the rule for (f(g(x))).
(g(x)=x+1)
(f(g(x))=f(x+1))
(=(x+1)^2)
(=x^2+2x+1)
(f(g(x))=x^2+2x+1)
When you apply the rule for one function to the rule of another function, you create a new function.
\answer{B. (f(g(x))=x^2+2x+1)}
\end{example}
\begin{tryit}{4}
Let (f(x)=2x-1) and (g(x)=3x). Find the value or identify the rule for the following functions.
a. (f(g(2)))
b. (f(g(x)))
\answer{a. (f(g(2))=11). b. (f(g(x))=6x-1).}
\end{tryit}
\begin{te-addendum}{4}{CommonError}
\textbf{Common Error. Try It! 4b}
When evaluating (f(3x)), students may write (2(3x-1)) rather than (2(3x)-1). Encourage them to write the function as (f(\square)=2(\square)-1) to emphasize that the expression that appears in the parentheses after the f should appear within parentheses on the right side of the equation in place of the variable x.
\end{te-addendum}
\begin{concept-box}
\textbf{CONCEPT Composite Function}
A composite function is the result of applying the rule for one function, f, to the rule of another function, g. The new rule is denoted as (f\circ g):
((f\circ g)(x)=f(g(x)))
The operation, (\circ), that forms a composite function is called composition of functions.
The domain of (f\circ g) is the set of all real numbers x in the domain of g such that (g(x)) is in the domain of f.
So the domain of the composition is the intersection of the domains of g and (f\circ g), but not f.
\end{concept-box}
\begin{te-addendum}{5}{PurposefulQuestions}
\textbf{Write a Rule for a Composite Function. Pose Purposeful Questions}
Q: What notation is introduced in the example and what does it mean?
\answer{The operation symbol (\circ) in (f\circ g) means "apply the rule for g, and then apply the rule for f."}
Q: What procedure do you perform to find the rule for (f\circ g)?
\answer{Substitute the expression (g(x)) for the variable in (f(x)), then apply the operations indicated by (f(x)).}
\end{te-addendum}
\begin{te-addendum}{5}{RtISupport}
\textbf{Support Struggling Students. USE WITH EXAMPLE 5}
Help students focus on specific aspects of the example by discussing the following questions.
Q: What does the expression (f(g(x))) in Part A mean?
\answer{Evaluate (g(x)) to get a number or an expression. Then evaluate the function f for that number or expression.}
Q: How is (g(f(x))) in Part B different from (f(g(x)))?
\answer{The difference is the order of using the functions. Evaluate the "inside" expression first, then use the result to evaluate the "outside" expression.}
\te{Source PDF page 69. Printed Part B reference to (g(f(x))); Example 5B displays (f(g(x))).}
\end{te-addendum}
\begin{example}{5}
\textbf{Write a Rule for a Composite Function}
What is the rule for the composition ((f\circ g)(x))?
A. (f(x)=\sqrt{x+7}) and (g(x)=2x-5)
((f\circ g)(x)=f(g(x)))
(=f(2x-5))
(=\sqrt{(2x-5)+7})
(=\sqrt{2x+2})
The rule for the composition ((f\circ g)(x)) is (\sqrt{2x+2}). The domain is (x\geq-1), which is the intersection of the domains of g and (f\circ g), which are all real numbers and (x\geq-1).
\te{MAKE SENSE AND PERSEVERE When finding the rule for the composition of functions, (f\circ g), the rule for g is substituted into the rule for f. How would you find the rule for (g\circ f)?}
\answer{A. (\sqrt{2x+2}); (x\geq-1).}
% Source page 70: 265 Example 5 continued
B. (f(x)=x^2+x+2) and (g(x)=4-x)
((f\circ g)(x)=f(g(x)))
(=f(4-x))
(=(4-x)^2+(4-x)+2)
(=16-8x+x^2+4-x+2)
(=x^2-9x+22)
The rule for the composition ((f\circ g)(x)) is (x^2-9x+22) and the domain is all real numbers.
\answer{B. (x^2-9x+22); all real numbers.}
\end{example}
\begin{tryit}{5}
Identify the rules for (f\circ g) and (g\circ f).
a. (f(x)=x^3), (g(x)=x+1)
b. (f(x)=3\sqrt{2x}), (g(x)=4x+11)
\answer{a. ((f\circ g)(x)=x^3+3x^2+3x+1) and ((g\circ f)(x)=x^3+1). b. ((f\circ g)(x)=3\sqrt{8x+22}) and ((g\circ f)(x)=12\sqrt{2x}+11).}
\end{tryit}
\begin{te-addendum}{5}{ElicitEvidence}
\textbf{Elicit and Use Evidence of Student Thinking}
Q: Part a: For the two given functions f and g, how can you tell when (f\circ g) will be negative?
\answer{(f(x)=x^3), so (f(g(x))) will be negative when (g(x)) is negative. Because (g(x)) is negative when (x+1<0), (f\circ g) is negative when (x<-1).}
\end{te-addendum}
\begin{te-addendum}{6}{PurposefulQuestions}
\textbf{Use a Composite Function Model. Pose Purposeful Questions}
Q: Why does order matter for discounts?
\answer{The amount of a percent coupon depends on the base price. To maximize a percent coupon, take it off the larger price; that is, before a dollars-off coupon.}
\end{te-addendum}
\begin{example}{6}
\textbf{Use a Composite Function Model. APPLICATION}
A clothing store posts discounts on social media. Followers are allowed to use multiple discounts. They simply need to tell the cashier in which order they want the discounts.
On their next trip to the store, in which order should they ask for the discounts to save the most money?
% FIG 70 1044 424 1144 608
\placeholder{illustration}{Black smartphone showing two offers, blue and green: EXCLUSIVELY FOR OUR SOCIAL MEDIA FOLLOWERS! 5 dollars OFF NEXT PURCHASE; EXCLUSIVELY FOR OUR SOCIAL MEDIA FOLLOWERS! 10 percent OFF NEXT PURCHASE. Phone status time 2:29PM.}
Write functions to model the discounts, letting x represent the price of the purchase.
(f(x)=x-5)
(g(x)=x-0.1x)
(=0.9x)
\te{Callout: price - 10 percent of price. STUDY TIP Using a model can help solve a problem. By writing functions to model the discounts, determine which sequence of functions provides a better deal for the customer.}
\begin{tabular}{ll}
10\% off, then \$5 off: Identify the rule for (f\circ g). & \$5 off, then 10\% off: Identify the rule for (g\circ f). \\
((f\circ g)(x)=f(g(x))) & ((f\circ g)(x)=g(f(x))) \\
(=f(0.9x)) & (=g(x-5)) \\
(=0.9x-5) & (=0.9(x-5)) \\
 & (=0.9x-4.5) \\
\end{tabular}
\te{Callout for left result: This is the better deal. Printed right-column label (f\circ g); likely (g\circ f).}
Suppose a shirt costs \$30.00 before the discounts. Applying the discounts as ((f\circ g)(x)) means the new cost is ((0.9)(30)-5=22), or \$22.00. Applying the discounts as ((f\circ g)(x)) means the new cost is ((0.9)(30)-4.5=22.5), or \$22.50.
The first method yields the better deal.
\answer{10\% off, then \$5 off.}
\te{Printed second composition label also (f\circ g); likely (g\circ f). Examples are available at SavvasRealize.com. Activity.}
\end{example}
\begin{tryit}{6}
As a member of the Games Shop rewards program, you get a 20\% discount on purchases. All sales are subject to a 6\% sales tax. Write functions to model the discount and the sales tax, then identify the rule for the composition function that calculates the final price you would pay at Games Shop.
\answer{Let x = purchase price. Discount: (d(x)=0.8x). Tax: (t(x)=1.06x). Actual price with discount: (t\circ d=0.848x).}
\end{tryit}
\begin{te-addendum}{6}{HabitsOfMind}
\textbf{HABITS OF MIND. Use with EXAMPLES 4--6}
Q: Reason Let (f(x)=\frac13(x-2)) and (g(x)=3x+2). Find (f\circ g) and (g\circ f). Are they equivalent? If so, does this prove that the composition of two functions is commutative?
\answer{(f(g(x))=x) and (g(f(x))=x), so in this case (f\circ g=g\circ f). One example is not a proof, so this does not prove that composition of functions is commutative.}
\end{te-addendum}
\begin{te-addendum}{6}{ELLAddendum}
\textbf{English Language Learners. Use with EXAMPLE 6}
LISTENING. BEGINNING. Point out the words order and sequence in the example and explain that they mean the same thing. Have students listen as you read the definition of order: "the arrangement of things following one after another, as in space or time."
Q: How is order important in the problem?
\answer{The final price changes depending on which discount is used first.}
Q: What is an example of when you can not change the order of operations in a problem?
\answer{When you simplify an expression, you must follow the same order of operations.}
SPEAKING. DEVELOPING. Choose some words from the example, such as discounts, deal, customers, cashier, etc. Have students take turns describing a word to the other students and have the others guess the word.
Q: How might a customer and a cashier interact?
\answer{A cashier helps a customer buy items in a store; he rings up the purchases, and the customer gives the cashier money for the purchases.}
Q: How are discounts and deals related?
\answer{A discount is money off the price of an item, and a deal is a good price on an item.}
WRITING. EXPANDING. There are some advanced verbs used in the example. Point out the words applying, posts, identify, and yields. Challenge students to write their own sentences correctly using these verbs.
Q: What does it mean to post on social media?
\answer{to put a statement, advertisement, or photo online to share with other people}
Q: What is a word that means the same thing as the word yields, as it is used in this example?
\answer{produces or gives}
\end{te-addendum}
% Source page 71: 266 Concept Summary
\begin{te-addendum}{ConceptSummary}{PurposefulQuestions}
\textbf{CONCEPT SUMMARY Function Operationss}
Q: How do operations on functions relate to what you have learned before?
\answer{Performing operations on functions is very similar to performing operations on integers, rational numbers, and polynomials.}
\te{Printed heading Operationss retained.}
\end{te-addendum}
\begin{te-addendum}{ConceptSummary}{CommonError}
\textbf{Common Error. Exercise 10}
Some students may believe that a zero of any function is not in the domain of a combination of functions. Encourage students to use a graphing utility to graph the combined function and look for asymptotes. The domain is only restricted for values where an asymptote occurs.
\te{Printed last sentence applies to this exercise; in general, canceled factors can also create holes excluded from the domain.}
\end{te-addendum}
\begin{concept-summary}
\textbf{CONCEPT SUMMARY Function Operations}
\begin{tabular}{llll}
 & Add or Subtract Functions & Multiply or Divide Functions & Compose Functions \\
ALGEBRA & ((f+g)(x)=f(x)+g(x)); ((f-g)(x)=f(x)-g(x)) & ((f\cdot g)(x)=f(x)\cdot g(x)); ((\frac fg)(x)=\frac{f(x)}{g(x)}) & ((f\circ g)(x)=f(g(x))); ((g\circ f)(x)=g(f(x))) \\
WORDS & The domain of the sum or difference of f and g is the intersection of the domain of f and the domain of g. & The domain is the set of all real numbers for which f and g and the new function are defined. & The domain of (f\circ g) is the set of all real numbers x, in the domain of g, such that (g(x)) is in the domain of f. \\
NUMBERS & Let (f(x)=3x+5) and (g(x)=x-3). & Let (f(x)=3x+5) and (g(x)=x-3). & Let (f(x)=3x+5) and (g(x)=x-3). \\
 & (f+g=(3x+5)+(x-3)=4x+2) & (f\cdot g=(3x+5)(x-3)=3x^2-4x-15) & (f\circ g=3(x-3)+5=3x-4) \\
 & (f-g=(3x+5)-(x-3)=2x+8) & (\frac fg=\frac{3x+5}{x-3}) for (x\neq3) & (g\circ f=(3x+5)-3=3x+2) \\
\end{tabular}
\textbf{Do You UNDERSTAND?}
1. ESSENTIAL QUESTION How do you combine, multiply, divide, and compose functions, and how do you find the domain of the resulting function?
\answer{To perform an arithmetic operation on a pair of functions, substitute their rules and perform the operation on them. The domain of the resulting function will be the set of all real numbers for which the original two functions and the resulting function are defined. To compose functions, apply one function to the rule of the other. The domain of the resulting function will be the set of all real numbers for which both the composite function and the inner function (the first function applied) are defined.}
2. Vocabulary In your own words, define and provide an example of a composite function.
\answer{A composite function is a method of combining functions such that the output of one function is the input of the other function. Example: Given (f(x)=2x) and (g(x)=x+1), the composite function (f\circ g=f(g(x))=2(x+1)=2x+2).}
3. Error Analysis Roble said the domain of (\frac fg) when (f(x)=5x^2) and (g(x)=x+3) is the set of real numbers. Explain why Roble is incorrect.
\answer{The domain is all real numbers except (x=-3), for which (g(x)) is equal to 0. (g(x)) is in the denominator, so when (x=-3), (\frac fg) is undefined.}
4. Use Structure Explain why changing the order in which two functions occur affects the result when subtracting and dividing the functions.
\answer{The operations of subtraction and division are not commutative.}
\textbf{Do You KNOW HOW?}
Let (f(x)=3x^2+5x+1) and (g(x)=2x-1).
5. Identify the rule for (f+g).
\answer{(3x^2+7x)}
6. Identify the rule for (f-g).
\answer{(3x^2+3x+2)}
7. Identify the rule for (g-f).
\answer{(-3x^2-3x-2)}
Let (f(x)=x^2+2x+1) and (g(x)=x-4).
8. Identify the rule for (f\cdot g).
\answer{(x^3-2x^2-7x-4)}
9. Identify the rule for (\frac fg), and state the domain.
\answer{(\frac{x^2+2x+1}{x-4}) for all real numbers x, such that (x\neq4)}
10. Identify the rule for (\frac gf), and state the domain.
\answer{(\frac{x-4}{x^2+2x+1}) or (\frac{x-4}{(x+1)^2}) for all real numbers x, when (x\neq-1)}
11. If (f(x)=2x^2+5) and (g(x)=-3x), what is (f(g(x)))?
\answer{(18x^2+5)}
\te{Concept Summary, Do You Understand, and Do You Know How are available at SavvasRealize.com. Presentation; Practice.}
\end{concept-summary}
% Source page 72: 267 Practice & Problem Solving
\begin{te-addendum}{Practice}{GeneralizeETP}
\textbf{STEP 3 Practice \& Problem Solving. Lesson Practice}
You may opt to have students complete the automatically scored Practice and Problem Solving items online powered by MathXL for School.
Choose from: Lesson Practice; Adaptive Practice.
You may also take advantage of the bank of exercises for assigning additional practice.
\textbf{Assignment Guide}
\begin{tabular}{ll}
On-Level & Advanced \\
13, 19--23, 33 & 12--36 \\
\end{tabular}
\textbf{Engage Through Student Choice}
Promote student agency by allowing students to choose practice items. You may structure this choice in many ways.
For example:
Assign each section a point value. Students choose at least one item from each section and items chosen should have a minimum of 20 total points.
\begin{tabular}{ll}
Understand, Apply & 2 points each \\
Practice & 1 point each \\
Assessment Practice & 1 point each \\
Performance Task & 3 points \\
\end{tabular}
\te{Practice \& Problem Solving and video tutorials are available at SavvasRealize.com. Scan the QR code to access a video tutorial. Practice.}
\end{te-addendum}
\begin{item-analysis}
\begin{tabular}{lll}
\toprule
Example & Items & DOK \\
\midrule
1 & 20--21 & 1 \\
1 & 19, 33 & 2 \\
2 & 22 & 2 \\
3 & 23 & 1 \\
3 & 13 & 2 \\
4 & 24--27, 35 & 1 \\
4 & 18, 31, 36 & 2 \\
5 & 28, 29 & 1 \\
5 & 12, 14, 34 & 2 \\
6 & 15, 17, 30 & 2 \\
6 & 16, 32 & 3 \\
\bottomrule
\end{tabular}
\end{item-analysis}
\begin{practice}{12}{2}
UNDERSTAND. Generalize Does (f\circ g) always equal (g\circ f)? Justify your response.
\answer{In general (f\circ g) and (g\circ f) are usually not equal. Sample: If (f(x)=x+1) and (g(x)=4x), then ((f\circ g)(x)=4x+1) and ((g\circ f)(x)=4x+4).}
\end{practice}
\begin{practice}{13}{2}
Construct Arguments Explain why the domain for the quotient of functions might not be the set of all real numbers.
\answer{because there are restrictions on the x-values in the denominator of a fraction such that it cannot be defined}
\end{practice}
\begin{practice}{14}{2}
Error Analysis Describe and correct the error a student made in finding the rule for the composition (f\circ g) of the functions (f(x)=3x^2-x+2) and (g(x)=2x+1).
(f\circ g=f(g(x)))
(=3(2x+1)^2-2x+1+2)
(=3(4x^2+4x+1)-2x+1+2)
(=12x^2+12x+3-2x+1+2)
(=12x^2+10x+6) cross
\answer{When (2x+1) was substituted into (3x^2-x+2) for x and when (-x) was replaced with (2x+1), the negative sign was not distributed.}
\end{practice}
\begin{practice}{15}{2}
Make Sense and Persevere Identify the rules for two functions, (f(x)) and (g(x)), for which (f\circ g=g\circ f).
\answer{Answers may vary. Sample: When (f(x)=2x) and (g(x)=5x), ((f\circ g)(x)=10x) and ((g\circ f)(x)=10x).}
\end{practice}
\begin{practice}{16}{3}
Higher Order Thinking Suppose two functions, (f(x)) and (g(x)) are only defined by the ordered pairs listed below.
(f=(6,7),(5,2),(4,1),(10,8))
(g=(5,4),(3,6),(1,5),(2,10))
Find the ordered pairs that comprise ((f\circ g)(x)).
\answer{((f\circ g)(x)=\{(5,1),(3,7),(1,2),(2,8)\})}
\end{practice}
\begin{practice}{17}{2}
Mathematical Connections Relate evaluating ((f\circ g)(3)) to finding the composition rule ((f\circ g)(x)). What are the benefits of each?
\answer{Evaluating function composition for a specific value is finding (g(3)) and then substituting (g(3)) for x in (f(x)). This approach is good when you only want to find a few values. Finding the rule of ((f\circ g)(x)) is similar in that you evaluate (f(x)) for the rule (g(x)), that is (f(g(x))). You can then recognize key features of the new function, including domain, range, and end behavior.}
\te{Answer printed on PDF page 73.}
\end{practice}
\begin{practice}{18}{2}
Make Sense and Persevere Identify rules for two functions (f(x)) and (g(x)), for which (f(g(x))=x).
\answer{Sample answer: If (f(x)=\frac12x) and (g(x)=2x), then (f(g(x))=x).}
\te{Answer printed on PDF page 73.}
\end{practice}
\begin{practice}{19}{2}
Construct Arguments Is it possible that (f(x)-g(x)=c), where c is a constant? If so, give an example. What must be true about two linear functions? If not, explain why it is not possible.
\answer{It is possible; Sample: If the graphs of the two linear functions are parallel, then their difference will be a horizontal line; (f(x)=3x-4) and (g(x)=3x+2), so (f(x)-g(x)=-6), which represents a horizontal line.}
\end{practice}
\begin{practice}{20}{1}
PRACTICE. Let (f(x)=\frac53x^2+2x-\frac58) and (g(x)=3x^2). Identify the rules for the following functions. SEE EXAMPLE 1
(f+g)
\answer{(f+g=\frac{14}3x^2+2x-\frac58)}
\end{practice}
\begin{practice}{21}{1}
Let (f(x)=\frac53x^2+2x-\frac58) and (g(x)=3x^2). Identify the rules for the following functions. SEE EXAMPLE 1
(f-g)
\answer{(f-g=-\frac43x^2+2x-\frac58)}
\end{practice}
\begin{practice}{22}{2}
Suppose the demand d, in units sold, for a company's jeans at price x, in dollars, is (d(x)=6,500-6.83x).
a. If revenue = price (\times) demand, write the rule for the function (R(x)), which represent the company's expected revenue in jean sales. Then state the domain of this function.
\answer{(R(x)=6,500x-6.83x^2) for (0\leq x\leq951.68)}
b. If the cost to manufacture the jeans is (c(x)=386+1.27x), find the equation for the company's profit. How much does the company earn if the price is \$79?
SEE EXAMPLE 2
\answer{b. (P(x)=-6.83x^2+6,498.73x-386); (P(79)=\$470,387.64)}
\te{Part b answer printed on PDF page 73. Domain endpoint in part a is rounded.}
\end{practice}
\begin{practice}{23}{1}
Identify the rule and domain for (\frac fg) when (f(x)=10x^2+3x-18) and (g(x)=2x+3). SEE EXAMPLE 3
\answer{((\frac fg)(x)=5x-6) for all real numbers x, such that (x\neq-\frac32)}
\end{practice}
\begin{practice}{24}{1}
Let (f(x)=4x-5) and (g(x)=-7x). Evaluate each expression. SEE EXAMPLE 4
(f(g(3)))
\answer{-89}
\end{practice}
\begin{practice}{25}{1}
Let (f(x)=4x-5) and (g(x)=-7x). Evaluate each expression. SEE EXAMPLE 4
(f(g(x)))
\answer{(-28x-5)}
\end{practice}
\begin{practice}{26}{1}
Let (f(x)=4x-5) and (g(x)=-7x). Evaluate each expression. SEE EXAMPLE 4
(g(f(2)))
\answer{-21}
\end{practice}
\begin{practice}{27}{1}
Let (f(x)=4x-5) and (g(x)=-7x). Evaluate each expression. SEE EXAMPLE 4
(g(f(x)))
\answer{(-28x+35)}
\end{practice}
\begin{practice}{28}{1}
Let (f(x)=x^2+x) and (g(x)=9-2x). Identify the rules for the following functions. SEE EXAMPLE 5
(f\circ g)
\answer{(4x^2-38x+90)}
\end{practice}
\begin{practice}{29}{1}
Let (f(x)=x^2+x) and (g(x)=9-2x). Identify the rules for the following functions. SEE EXAMPLE 5
(g\circ f)
\answer{(-2x^2-2x+9)}
\end{practice}
\begin{practice}{30}{2}
A sporting goods store is running a summer sale on its snowboards. Kadia is interested in a snowboard that normally costs \$400. The store is offering two special offers.
In which order should these special offers be applied to the cost of the snowboard in order to benefit Kadia? Explain. SEE EXAMPLE 6
% FIG 72 997 772 1116 1038
\placeholder{illustration}{Snowboards in a store display, central blue/white board labeled 58. Circular offers at bottom: 50 dollars INSTANT REBATE; 10 percent DISCOUNT.}
\answer{It is better for Kadia if the 10\% discount is applied before the \$50 instant rebate because he will save \$5 more.}
\te{Answer printed on PDF page 73.}
\end{practice}
% Source page 73: 268 Practice & Problem Solving
\begin{te-addendum}{Practice}{GeneralizeETP}
\te{Practice \& Problem Solving is available at SavvasRealize.com. Answers for 17--18, 22--23, and 30 are transcribed with their items.}
\end{te-addendum}
\begin{practice}{31}{2}
APPLY. Model With Mathematics The cost (in dollars) to produce x shovels in a factory is given by the function (C(x)=20x+500). The factory produces 30 shovels in one hour.
a. Find the rule for (C(h(x))), where the function (h(x)) is the number of shovels produced in x hours.
\answer{a. (C(h(x))=600x+500)}
b. Find the cost when (h=8) hours.
\answer{b. \$5,300}
c. Explain what the function (C(h(x))) represents.
\answer{c. (C(h(x))) represents the cost to produce x hours' worth of shovels.}
\te{Printed h=8 hours; likely x=8 hours.}
\end{practice}
\begin{practice}{32}{3}
Use Structure A music store is running the following promotions.
% FIG 73 518 477 774 670
\placeholder{illustration}{O'RILEY'S MUSIC drum-set advertisement: HUGE SAVINGS AT O'RILEY'S MUSIC. 5 dollars off a purchase of 20 dollars or more. Save an additional 15 percent off your total purchase when you open a charge account.}
a. Use composition of functions to find the sale price of a \$90 purchase when the \$5 off discount is applied prior to the 15\% off discount.
\answer{a. \$72.25}
b. Use composition of functions to find the sale price of a \$90 purchase when the 15\% off discount is applied prior to the \$5 off discount.
\answer{b. \$71.50}
c. In which order is the deal better for the consumer? Explain.
\answer{c. Since (\$71.50<\$72.25), it is a better deal for the consumer if the 15\% discount is applied prior to the \$5 off discount.}
\end{practice}
\begin{practice}{33}{2}
Reason From 2000 to 2015, the number of births, b, (in the hundreds) in Fairfield County can be modeled by the function (b(x)=300-5x). The number of deaths, d, (in the hundreds) can be modeled by the function (d(x)=10x+5). The variable x represents the number of years since 2000.
a. Which function operation can be used to represent the net increase in the population?
\answer{a. subtraction}
b. Write and simplify a function which represents the net increase in the population, p, against x, the number of years since 2000. State the domain of this function.
\answer{b. (p(x)=b(x)-d(x)=-15x+295); the domain is the set of all real numbers x such that (0\leq x\leq15).}
\end{practice}
\begin{practice}{34}{2}
ASSESSMENT PRACTICE. Given that (f(x)=x^2+8x+3) and (g(x)=-x-7), which of the following are true? Select all that apply.
\item (f+g=x^2+7x-4)
\item (f(g(x))=x^2+6x-4)
\item The domain of (\frac fg) is the set of all real numbers.
\item (f(x)\cdot g(x)=-x^3-15x^2+53x+21)
\item In the composition (g\circ f), the output (f(x)) is used as the input for g.
\answer{First, second, and fifth statements.}
\end{practice}
\begin{practice}{35}{1}
SAT/ACT Find the value of (f(g(5))) if (f(x)=4x+1) and (g(x)=x^2+6).
A. 101
B. 124
C. 125
D. 676
E. 682
\answer{C}
\end{practice}
\begin{practice}{36}{2}
Performance Task The temperature in degrees Celsius is 32 less than the Fahrenheit temperature, multiplied by five-ninths. The temperature in degrees Kelvin is the number of degrees Celsius plus 273.
% FIG 73 824 608 1069 792
\placeholder{diagram}{Three vertical thermometers labeled Celsius, degrees C; Fahrenheit, degrees F; Kelvin, K. Orange horizontal reference line aligns 0 degrees C, 32 degrees F and 273 K. Celsius scale labels -40,-20,0 degrees C,20,40,60, small ticks every 5; Fahrenheit labels -50,-25,0,25,50,75,100,125,150, small ticks every 5; Kelvin labels 230,250,270,290,310,330, small ticks every 5. Black columns inside pale blue thermometer backgrounds.}
Part A Derive a conversion formula for finding the number of degrees Kelvin, given the temperature in Fahrenheit.
\answer{Part A (K=\frac59(F-32)+273)}
Part B Using your conversion formula from part (a), find the temperature in degrees Kelvin when the temperature is 27 degrees F. Round to the nearest whole number if necessary.
\answer{Part B 270 K}
\end{practice}
% Source page 74: 268A Assess & Differentiate
\begin{te-addendum}{Assess}{RtISupport}
\textbf{STEP 4 Assess & Differentiate. LESSON QUIZ}
Use the Lesson Quiz to assess students' understanding of the mathematics in the lesson.
Students can take the Lesson Quiz online or you can download a printable copy from SavvasRealize.com. The Lesson Quiz is also available in the Assessment Sourcebook.
\textbf{Item Analysis}
\begin{tabular}{lll}
Item & DOK & Skills Review and Practice \\
1 & 3 & F52, F26 \\
2 & 3 & F53, F26 \\
3A & 3 & F53, F26 \\
3B & 2 & F53, F26 \\
4 & 3 & F53, F26 \\
5 & 2 & F64, F26 \\
\end{tabular}
Use the student scores on the Lesson Quiz to prescribe differentiated assignments.
If students take the Lesson Quiz online, it will be automatically scored and appropriate differentiated practice will be assigned based on student performance.
\begin{tabular}{lll}
Intervention & 0--3 points & Reteach to Build Understanding; Mathematical Literacy and Vocabulary; Lesson Virtual Nerd videos \\
On-Level & 4 points & Enrichment \\
Advanced & 5 points & Enrichment \\
\end{tabular}
\end{te-addendum}
\begin{te-addendum}{Assess}{GeneralizeETP}
\textbf{ASSESSMENT RESOURCES. Lesson Quiz is available at SavvasRealize.com.}
\textbf{5-5 Lesson Quiz: Function Operations}
1. Let (f(x)=3x^2-2x+6) and (g(x)=7x-4). Identify the rule for (g(f(x))+g(x)).
A. (21x^2-7x+34)
B. (21x^2-21x+50)
C. (147x^2-182x+62)
D. (21x^3-26x^2+57x-28)
\answer{1. A}
2. The demand d for a company's product at cost x is predicted by the function (d(x)=500-2x). The price p in dollars that the company can charge for the product is given by (p(x)=x+5). Use the formula Revenue = price (\times) demand to find the revenue function for the product.
(R(x)=\underline{\hspace{1cm}}+\underline{\hspace{1cm}}x-\underline{\hspace{1cm}}x^2)
\answer{2. 2,500; 490; 2.}
3. Part A Identify the rule for (\frac fg) when (f(x)=-3x-6) and (g(x)=x^2-x-6).
A. (\frac{f(x)}{g(x)}=\frac{-2}{(x+2)})
B. (\frac{f(x)}{g(x)}=\frac{-3(x-2)}{(x-6)(x+1)})
C. (\frac{f(x)}{g(x)}=\frac3{(x+2)})
D. (\frac{f(x)}{g(x)}=\frac{-3}{x-3})
\answer{3. Part A D}
Part B Select all values that are NOT in the domain of (\frac fg).
A. -1
B. 3
C. -3
D. 6
E. -2
\answer{Part B B, E}
4. Let (h(x)=\sqrt{x+3}) and (k(x)=2x+7). Find the value (\frac{h(k(3))-k(h(1))}{k(h(-2))-h(k(-3))}).
A. (\frac{2\sqrt6-2\sqrt3}{\sqrt6-7})
B. 4
C. (2\sqrt6-\sqrt{12}+7)
D. -1
\answer{4. D}
5. Let (f(x)=3x^2+2) and (g(x)=\sqrt{x-4}). What is the rule for the composition (f\circ g)?
A. ((f\circ g)(x)=\sqrt{3x^2-2}), domain is all real numbers
B. ((f\circ g)(x)=\sqrt{3x^2-2}), domain is (x\geq4)
C. ((f\circ g)(x)=3x-10), domain is all real numbers
D. ((f\circ g)(x)=3x-10), domain is (x\geq4)
\answer{5. D}
\end{te-addendum}
% Source page 75: 268B Differentiated Resources
\begin{te-addendum}{Assess}{RtISupport}
\textbf{STEP 4 Assess & Differentiate. DIFFERENTIATED RESOURCES}
I = Intervention. O = On-Level. A = Advanced. SavvasRealize.com. Practice.
\textbf{Reteach to Build Understanding. I}
Provides scaffolded reteaching for the key lesson concepts. MathXL for School.
\textbf{5-5 Reteach to Build Understanding: Function Operations}
1. Complete the table to practice each operation.
\begin{tabular}{lllll}
 & Step 1: Rule & Step 2: Equations & Step 3: Set Up & Step 4: Answer \\
Add & (f(x)+g(x)) & (f(x)=x-11); (g(x)=2x+11) & (x-11+2x+11) & (=\answer{3}x) \\
Subtract & (f(x)-g(x)) & (f(x)=4x+8); (g(x)=2x+5) & (4x+8-(2x+5)); (4x+8-2x-5) & (=\answer{2}x+\answer{3}) \\
Multiply & (f(x)\cdot g(x)) & (f(x)=x-3); (g(x)=2x+2) & ((x-3)(2x+2)); (2x^2+2x-6x-6) & (=\answer{2}x^2-\answer{4}x-\answer{6}) \\
Divide & (\frac{f(x)}{g(x)}) & (f(x)=x^2+6x+8); (g(x)=x+4) & (\frac{x^2+6x+8}{x+4}) & (=x+\answer{2}) \\
Compose & (f(g(x))) & (f(x)=2x+3); (g(x)=3x+2) & (2(3x+2)+3); (6x+4+3) & (=\answer{6}x+\answer{7}) \\
\end{tabular}
2. Fatima was asked to write the rule for (f\circ g). She set it up as (f(x)\cdot g(x)). What error did she make? Choose the correct answer.
a. She set up the rule for multiplying the two equations.
b. She set up the rule for adding the two equations.
c. She set up the rule for dividing the two equations.
d. She set up the rule for subtracting the two equations.
\answer{a}
How should she have set it up? Choose the correct answer.
a. When composing them, she should have set them up as (\frac{f(x)}{g(x)}).
b. When composing them, she should have set them up as (f(x)-g(x)).
c. When composing them, she should have set them up as (f(g(x))).
d. When composing them, she should have set them up as (f(x)+g(x)).
\answer{c}
3. What is the rule for the composition (f\circ g)?
(f(x)=\sqrt{x-3}) and (g(x)=2x+4)
(f\circ g=f(g(x)))
(=f(2x+\underline{\hspace{0.4cm}}))
(=\sqrt{(2x+4)-\underline{\hspace{0.4cm}}})
(=\sqrt{\underline{\hspace{0.4cm}}x+\underline{\hspace{0.4cm}}}) Domain: (x\geq-\frac12)
\answer{4; 3; 2; 1.}
\end{te-addendum}
\begin{te-addendum}{Assess}{GeneralizeETP}
\textbf{Additional Practice. I O}
Provides extra practice for each lesson. MathXL for School.
\textbf{5-5 Additional Practice: Function Operations}
Let (f(x)=3x^2-9x-11) and (g(x)=7-4x). Identify rules for the following functions.
1. (f+g) \answer{(f(x)+g(x)=3x^2-13x-4)}
2. (f-g) \answer{(f(x)-g(x)=3x^2-5x-18)}
3. Suppose demand d for a company's product at cost x is predicted by the function (d(x)=0.36x^2+810), and that the price p that the company can charge for the product is given by (p(x)=x+14). Find the company's revenue function.
\answer{(f(x)=0.36x^2-x+796)}
\te{Printed answer subtracts p from d; revenue would normally use the product, (0.36x^2+810)(x+14).}
4. Identify the rule and domain for (\frac fg) when (f(x)=x^2-5x-36) and (g(x)=x-9).
\answer{((\frac fg)(x)=x+4); all real numbers excluding 9}
Let (f(x)=3x-2) and (g(x)=5x). Identify the rule for the following functions.
5. (f(g(3))) \answer{43}
6. (f(g(x))) \answer{(15x-2)}
7. Identify the rules for (f\circ g) and (g\circ f) when (f(x)=2x^3) and (g(x)=x-1).
\answer{((f\circ g)(x)=2(x-1)^3); ((g\circ f)(x)=2x^3-1)}
8. As a member of the Games Shop rewards program, you get a 12\% discount on purchases. All sales are subject to an 8\% sales tax. Write functions to model the discount and the sales tax, then identify the rule for the composition function that calculates the final price you pay Games Shop.
\answer{(D(x)=x-0.12x); (T(x)=0.08(x-0.12x)); (P=0.9504x)}
9. Describe and correct the error a student made in finding the rule for the composition (f\circ g) when (f(x)=2x^2-3x+1) and (g(x)=2x-1).
((f\circ g)(x)=f(g(x)))
(=2(2x-1)^2-3x+1)
(=2(4x^2-4x+1)-3x+1)
(=8x^2-11x+3)
\answer{The student did not replace the x is -3x with (2x-1). (=2(2x-1)^2-3(2x-1)+1); (=2(4x^2-4x+1)-6x+3+1); (=8x^2-14x+6)}
\te{Printed x is -3x; likely x in -3x.}
The cost in dollars to produce x shovels in a factory is given by the function (C(x)=23x+480). The number of shovels that can be produced in h hours is given by the function (N(h)=30h).
10. Find the rule for (C(N(h))). \answer{(690h+480)}
11. Find the cost when (h=8) hours. \answer{\$6,000}
Let (f(x)=3x^2+2x-3) and (g(x)=2x+4). Identify the rules for the following functions.
12. (f+g) \answer{(3x^2+4x+1)}
13. (f-g) \answer{(3x^2-7)}
\end{te-addendum}
\begin{te-addendum}{Assess}{RtIExtend}
\textbf{Enrichment. O A}
Presents engaging problems and activities that extend the lesson concepts. MathXL for School.
\textbf{5-5 Enrichment: Function Operations}
Use compositions of the following functions to fill cells 1--6.
(f(x)=x^2-1)
(g(x)=x+1)
(h(x)=x-1)
(k(x)=x^2+1)
Rules of the Game:
1. You can fill each cell only with one composition of the functions.
2. You can use a function in more than one composition.
3. You cannot use the same composition more than once.
4. The sum of terms in each column and row with more than one cell is equal to (3x^2).
5. Function (f(x)=x^2-1) must be used in all compositions of the left column.
% FIG 75 911 574 1048 710
\placeholder{diagram}{Seven equal rectangular cells forming two offset vertical columns joined by one horizontal row: left column top-to-bottom cells 1,2,3; middle connecting cell 4 right of 3; right column top-to-bottom 5,6,7, with 5 right of 4. Black numbers in upper left. Cell 7 prints x squared minus 2x. Magenta keys: 1 f composed with g; 2 f composed with h; 3 g composed with f; 4 g composed with k; 5 h composed with f; 6 k composed with g.}
\answer{1. (f\circ g); 2. (f\circ h); 3. (g\circ f); 4. (g\circ k); 5. (h\circ f); 6. (k\circ g).}
\end{te-addendum}
\begin{te-addendum}{Assess}{ELLAddendum}
\textbf{Mathematical Literacy and Vocabulary. I O}
Helps students develop and reinforce understanding of key terms and concepts.
\textbf{5-5 Mathematical Literacy and Vocabulary: Function Operations}
Ashton wanted to find a function rule for (f\circ g), given (f(x)=x^2-1) and (g(x)=x-2). He wrote these steps on note cards, but they got mixed up.
% FIG 75 226 1028 371 1089
\placeholder{diagram}{Six overlapping note cards. Upper left: =(x-2) squared -1. Upper middle: =x squared -4x+4-1. Middle left: =f composed with g=f(g(x)). Middle right: =x squared -2(2)x+2 squared -1. Lower middle: =x squared -4x+3. Lower right: =f(x-2).}
Organize the cards to complete the steps Ashton used to find the function rule.
1. First, \answer{(f\circ g=f(g(x)))}
2. Second, \answer{(=f(x-2))}
3. Next, \answer{(=(x-2)^2-1)}
4. Next, \answer{(=x^2-2(2)x+2^2-1)}
5. Then, \answer{(=x^2-4x+4-1)}
6. Finally, \answer{(=x^2-4x+3)}
\end{te-addendum}
\begin{te-addendum}{Assess}{GeneralizeETP}
\textbf{Digital Differentiated Resources}
The Reteach to Build Understanding, Additional Practice, and Enrichment activities are available as digital assignments. These activities are automatically assigned when students complete the lesson quiz online and are automatically scored.
% FIG 75 586 979 1034 1298
\placeholder{illustration}{Tablet displaying online graphing exercise: Solve the system of equations by graphing, (y=3x+4), (y=-2x+4). Use the graphing tool to graph the system. Empty coordinate grid labeled x and y, both axes from -10 to 10 with unit grid spacing and labels every 2; arrow on positive y; part of upper right grid hidden by Question Help menu. Menu: Help Me Solve This, View an Example, Video, Textbook, Glossary, Math Tools, Print. Click to enlarge graph. Click the graph, choose a tool in the palette and follow the instructions to create your graph. Controls: Exit, 1 part remaining, Clear All, Check Answer, Review progress, Question 5 of 14, Go, Back, Next.}
\end{te-addendum}

\end{document}
